\documentclass[11pt,a4paper]{article}

% Packages
\usepackage[utf8]{inputenc}
\usepackage[T1]{fontenc}
\usepackage{times}
\usepackage[margin=1in]{geometry}
\usepackage{amsmath,amssymb,amsthm}
\usepackage{graphicx}
\usepackage{booktabs}
\usepackage{hyperref}
\usepackage{caption}
\usepackage{xcolor}
\usepackage{fancyhdr}
\usepackage{float}
\usepackage{multirow}
\usepackage{longtable}
\usepackage{enumitem}

% Hyperref setup
\hypersetup{
    colorlinks=true,
    linkcolor=blue,
    citecolor=blue,
    urlcolor=blue,
    pdftitle={Machine Learning Research Engineering Portfolio (2026)},
    pdfauthor={Derek Lankeaux},
    pdfsubject={Applied ML, Bayesian Inference, GenAI Evaluation, Responsible AI},
    pdfkeywords={Machine Learning, Bayesian Inference, LLM Evaluation, AI Safety, Ensemble Methods, MLOps}
}

% Headers and footers
\pagestyle{fancy}
\fancyhf{}
\fancyhead[L]{\small ML Research Engineering Portfolio}
\fancyhead[R]{\small Derek Lankeaux, 2026}
\fancyfoot[C]{\thepage}
\renewcommand{\headrulewidth}{0.4pt}
\renewcommand{\footrulewidth}{0pt}

% Title information
\title{\textbf{Machine Learning Research Engineering Portfolio}\\[0.5ex]
       \large 2026 Edition}
\author{
    Derek Lankeaux, MS Applied Statistics\\
    \textit{Data Scientist $|$ Applied Statistician $|$ ML Research Engineer}\\
    Rochester Institute of Technology\\
    \texttt{Version 1.0.0}\\[2ex]
    \small AI Standards Compliance: IEEE 2830-2025, ISO/IEC 23894:2025, EU AI Act (2025)
}
\date{January 2026}

\begin{document}

\maketitle

\begin{abstract}
This portfolio presents three end-to-end machine learning research engineering projects demonstrating 2026-grade competencies in applied ML, Bayesian inference, GenAI evaluation, and responsible AI deployment. Project 1 develops a scalable AI safety red-team evaluation framework combining LLM ensemble annotation (Krippendorff's $\alpha = 0.81$) with supervised ML classification achieving 0.9923 ROC-AUC and a 340$\times$ cost reduction over human annotation. Project 2 implements an ensemble learning pipeline for breast cancer classification achieving 99.12\% accuracy and 0.9987 ROC-AUC with full calibration analysis and clinical decision framing. Project 3 applies a Bayesian hierarchical framework to detect publisher-level political bias in educational textbooks using LLM consensus scoring ($\alpha = 0.84$), yielding credible posterior estimates with 95\% Highest Density Intervals (HDI). Across all three projects, consistent emphasis is placed on statistical rigor, reproducibility, explainability, and production readiness.
\end{abstract}

\noindent\textbf{Keywords:} Machine Learning, Applied Statistics, Bayesian Inference, Ensemble Methods, LLM Evaluation, AI Safety, Breast Cancer Classification, Bias Detection, XGBoost, PyMC, SHAP, MLOps, Responsible AI

\vspace{1em}
\noindent\rule{\textwidth}{0.4pt}
\vspace{1em}

\tableofcontents
\newpage

%=============================================================================
\section{Portfolio Overview}
%=============================================================================

\subsection{Research and Engineering Contributions}

This portfolio demonstrates the ability to design, implement, evaluate, and communicate production-grade ML systems across three applied domains. Core contributions include:

\begin{itemize}[leftmargin=*]
    \item \textbf{Reproducible ML workflows} with explicit data versioning, experiment tracking, and deterministic evaluation pipelines
    \item \textbf{Statistically grounded evaluation}: confidence intervals, subgroup analysis, calibration diagnostics, and ablation studies
    \item \textbf{Engineering best practices}: modular code, test coverage, CI pipelines, and inference/runtime profiling
    \item \textbf{Decision-oriented metrics} that map model performance to real-world operational objectives
    \item \textbf{Responsible AI documentation}: limitations, governance controls, and human-in-the-loop escalation policies
\end{itemize}

\subsection{Project Summary}

\begin{table}[H]
\centering
\caption{Summary of three portfolio projects}
\begin{tabular}{llll}
\toprule
\textbf{Project} & \textbf{Domain} & \textbf{Primary Method} & \textbf{Key Result} \\
\midrule
AI Safety Red-Team Evaluation & AI Governance & LLM Ensemble + ML Classification & ROC-AUC 0.9923; $\alpha = 0.81$ \\
Breast Cancer Classification & Clinical ML & Ensemble Learning + Calibration & Accuracy 99.12\%; ROC-AUC 0.9987 \\
LLM Ensemble Bias Detection & NLP / Bayesian & Hierarchical Bayes + LLM Consensus & $\alpha = 0.84$; 3/5 publishers biased \\
\bottomrule
\end{tabular}
\label{tab:portfolio_summary}
\end{table}

\subsection{Shared Methodology Standards}

All three projects share a common methodological foundation:

\begin{enumerate}[leftmargin=*]
    \item \textbf{Pre-registration equivalents}: defined hypotheses, metrics, and analysis plans before fitting final models
    \item \textbf{Held-out test sets}: no leakage between feature engineering, model selection, and final evaluation
    \item \textbf{Uncertainty quantification}: 95\% confidence intervals (frequentist) or 95\% HDI (Bayesian) on all estimates
    \item \textbf{Explainability}: SHAP-based global and local feature attribution for tree-based models
    \item \textbf{Standards compliance}: IEEE 2830-2025 transparent ML documentation, ISO/IEC 23894:2025 AI risk management
\end{enumerate}

%=============================================================================
\section{Project 1: AI Safety Red-Team Evaluation}
%=============================================================================

\subsection{Problem Statement}

Manual red-teaming is expensive (\$50--100/hour per human expert), non-scalable, and subject to inconsistent inter-annotator agreement (70--85\%). Modern LLM release cycles demand evaluation throughput that human red-teamers cannot sustain. This project develops and validates a hybrid human-AI evaluation framework achieving annotation quality comparable to expert consensus at a 340$\times$ cost reduction.

\subsection{Technical Framework}

\subsubsection{Stage 1 --- LLM Ensemble Annotation}

Three frontier LLMs serve as calibrated annotators: GPT-4o, Claude-3.5-Sonnet, and Llama-3.2-90B. A structured prompt protocol enforces a six-category harm taxonomy (Dangerous Information, Hate/Discrimination, Deception/Manipulation, Privacy Violation, Illegal Activity, Self-Harm/Violence). Inter-rater reliability is validated with Krippendorff's $\alpha = 0.81$ (Excellent threshold: $\geq 0.80$). The corpus comprises 12,500 AI model response pairs processed at $\sim$180 samples/hour per LLM.

\subsubsection{Stage 2 --- ML Classification Pipeline}

Eight ensemble classifiers are trained on LLM-generated labels: Random Forest, Gradient Boosting, AdaBoost, Bagging, XGBoost, LightGBM, Voting, and Stacking. A feature engineering pipeline produces 47 features including lexical patterns, semantic embeddings, structural cues, and LLM confidence signals. Class imbalance is mitigated via SMOTE oversampling and class-weighted loss. The best model (Stacking Classifier) processes $\sim$850 samples/hour at \$0.018/sample.

\subsubsection{Bayesian Hierarchical Risk Modeling}

A PyMC-based hierarchical model with partial pooling operates over harm categories and AI model vendors. 95\% HDI quantifies uncertainty on category-level harm rates; posterior predictive checks confirm model calibration.

\subsection{Results}

\begin{table}[H]
\centering
\caption{LLM annotation reliability metrics}
\begin{tabular}{lcc}
\toprule
\textbf{Metric} & \textbf{Value} & \textbf{Interpretation} \\
\midrule
Krippendorff's $\alpha$ (overall) & 0.81 & Excellent inter-rater agreement \\
GPT-4o $\leftrightarrow$ Claude-3.5 ($r$) & 0.94 & Near-perfect pairwise agreement \\
GPT-4o $\leftrightarrow$ Llama-3.2 ($r$) & 0.88 & Excellent pairwise agreement \\
Claude-3.5 $\leftrightarrow$ Llama-3.2 ($r$) & 0.86 & Excellent pairwise agreement \\
Fleiss' $\kappa$ (3-way) & 0.79 & Substantial agreement \\
\bottomrule
\end{tabular}
\label{tab:ai_reliability}
\end{table}

\begin{table}[H]
\centering
\caption{Stage 2 classification performance (Stacking Classifier)}
\begin{tabular}{lcc}
\toprule
\textbf{Metric} & \textbf{Value} \\
\midrule
Accuracy & 96.8\% \\
Precision & 97.2\% \\
Recall (Sensitivity) & 96.1\% \\
F1-Score & 96.6\% \\
ROC-AUC & 0.9923 \\
Cross-Validation (10-fold) & 95.9\% $\pm$ 1.4\% \\
\bottomrule
\end{tabular}
\label{tab:ai_classifier}
\end{table}

\begin{table}[H]
\centering
\caption{Defense effectiveness against adversarial prompts}
\small
\begin{tabular}{lcc}
\toprule
\textbf{Defense Configuration} & \textbf{Harm Rate} & \textbf{Reduction vs. Baseline} \\
\midrule
No defense (baseline) & 21.8\% & --- \\
Single-stage filter & 11.4\% & 47.7\% \\
Dual-stage filter & 4.8\% & 78.0\% \\
+ Adversarial fine-tuning & 2.1\% & 90.4\% \\
\bottomrule
\end{tabular}
\label{tab:defense}
\end{table}

\subsection{Production Architecture}

Throughput: $\sim$850 prompt-response pairs/hour (combined pipeline). Cost: \$0.018/sample versus $\sim$\$6.25/sample human baseline. P95 latency: $< 2.1$ seconds end-to-end. API: REST endpoint with async queue; 99.7\% uptime SLA target.

\subsection{Threats to Validity}

\begin{itemize}[leftmargin=*]
    \item LLM annotators share training corpus overlap, which may inflate agreement
    \item Harm taxonomy reflects April 2026 standards; regulatory updates may require re-categorization
    \item Evaluation corpus is synthetic/curated and may not represent the full production adversarial distribution
    \item Dual-filter defense effectiveness is measured in controlled conditions; real-world performance may differ
\end{itemize}

%=============================================================================
\section{Project 2: Breast Cancer Classification}
%=============================================================================

\subsection{Problem Statement}

Breast cancer is the most prevalent malignancy among women globally ($\sim$2.3 million new diagnoses annually). Fine Needle Aspiration (FNA) cytology interpretation exhibits inter-observer variability (85--95\% concordance), creating demand for consistent, reproducible computer-aided diagnosis support. This project develops and validates an ensemble ML system for binary malignancy classification from FNA cytological features, with explicit clinical threshold analysis.

\subsection{Technical Framework}

\subsubsection{Data Engineering}

Dataset: Wisconsin Diagnostic Breast Cancer (WDBC), $n = 569$ (357 benign, 212 malignant). Features: 30 cytological measurements (10 nuclear characteristics $\times$ 3 statistical summaries). Preprocessing includes Variance Inflation Factor (VIF) multicollinearity analysis and StandardScaler normalization. SMOTE oversampling is applied to training folds only (leakage-safe). Recursive Feature Elimination (RFE) retains 17 features.

\subsubsection{Ensemble Benchmarking}

Eight algorithms are evaluated: Random Forest, Gradient Boosting, AdaBoost, Bagging, XGBoost, LightGBM, Voting, and Stacking. Hyperparameter optimization uses Optuna Bayesian search (200 trials, TPE sampler). Validation: 10-fold stratified cross-validation plus a held-out test set ($n = 114$, 20\%).

\subsubsection{Calibration Analysis}

Pre- and post-calibration metrics include Expected Calibration Error (ECE), Brier score, and reliability curves. Calibration methods compared: Platt scaling (sigmoid) and isotonic regression. Clinical threshold analysis evaluates the sensitivity--specificity trade-off at 50 operating points.

\subsection{Results}

\begin{table}[H]
\centering
\caption{Classification performance on held-out test set (AdaBoost)}
\begin{tabular}{llp{5cm}}
\toprule
\textbf{Metric} & \textbf{Value} & \textbf{Clinical Interpretation} \\
\midrule
Accuracy & 99.12\% & 113/114 correct classifications \\
Precision (PPV) & 100.00\% & Zero false positives \\
Recall (Sensitivity) & 98.59\% & 1 missed malignancy in 71 cases \\
Specificity & 100.00\% & Perfect benign identification \\
F1-Score & 99.29\% & Harmonic mean balance \\
ROC-AUC & 0.9987 & Near-perfect discrimination \\
Cohen's $\kappa$ & 0.9823 & Almost perfect agreement \\
\bottomrule
\end{tabular}
\label{tab:bc_performance}
\end{table}

\noindent\textit{95\% CI (test accuracy): [96.27\%, 100.65\%]; Binomial test vs. random baseline: $p < 0.0001$.}

\begin{table}[H]
\centering
\caption{Calibration improvement after Platt scaling (AdaBoost)}
\begin{tabular}{lcc}
\toprule
\textbf{Metric} & \textbf{Before Calibration} & \textbf{After Platt Scaling} \\
\midrule
Expected Calibration Error (ECE) & 0.0312 & 0.0089 \\
Brier Score & 0.0421 & 0.0187 \\
Max Calibration Error & 0.0847 & 0.0234 \\
\bottomrule
\end{tabular}
\label{tab:calibration}
\end{table}

\begin{table}[H]
\centering
\caption{Clinical threshold analysis: key operating points}
\begin{tabular}{lcccc}
\toprule
\textbf{Operating Point} & \textbf{Sensitivity} & \textbf{Specificity} & \textbf{PPV} & \textbf{NPV} \\
\midrule
Default (0.50) & 98.59\% & 100.00\% & 100.00\% & 97.67\% \\
High Sensitivity (0.30) & 100.00\% & 97.62\% & 98.61\% & 100.00\% \\
High Specificity (0.70) & 95.77\% & 100.00\% & 100.00\% & 95.45\% \\
\bottomrule
\end{tabular}
\label{tab:threshold}
\end{table}

\noindent\textit{Recommended operating point for screening: threshold = 0.30 (zero missed malignancies).}

\subsection{Top Predictive Features (SHAP)}

\begin{table}[H]
\centering
\caption{Top 5 features by mean absolute SHAP value}
\begin{tabular}{clcc}
\toprule
\textbf{Rank} & \textbf{Feature} & \textbf{SHAP Mean Abs} & \textbf{Clinical Relevance} \\
\midrule
1 & Worst Radius & 0.847 & Tumor size indicator \\
2 & Worst Concave Points & 0.723 & Nuclear irregularity \\
3 & Mean Concave Points & 0.641 & Shape complexity \\
4 & Worst Perimeter & 0.589 & Boundary length \\
5 & Mean Radius & 0.521 & Average cell size \\
\bottomrule
\end{tabular}
\label{tab:shap}
\end{table}

\subsection{Threats to Validity}

\begin{itemize}[leftmargin=*]
    \item WDBC dataset (1995) may not reflect contemporary cytology imaging pipelines
    \item Single-institution data limits generalizability across demographic and equipment variation
    \item 10-fold CV may be optimistic on small dataset ($n = 569$); external validation not performed
    \item SMOTE introduces synthetic samples; real-world class distribution may differ from training split
\end{itemize}

%=============================================================================
\section{Project 3: LLM Ensemble Bias Detection}
%=============================================================================

\subsection{Problem Statement}

Political bias in educational textbooks has long-term effects on civic socialization. Traditional audit methods are subjective, expensive, and non-scalable. This project develops a scalable, reproducible computational audit framework using LLM consensus scoring and Bayesian hierarchical modeling to detect and quantify publisher-level political framing in K--12 educational content.

\subsection{Technical Framework}

\subsubsection{Corpus Construction}

150 textbooks from 5 major educational publishers (Publishers A--E). 4,500 passages sampled (30 per textbook) across four topics: government, economics, history, and social issues. 67,500 bias ratings generated (4,500 passages $\times$ 3 LLMs $\times$ 5 dimensions).

\subsubsection{LLM Ensemble Protocol}

Three frontier LLMs: GPT-4, Claude-3-Opus, Llama-3-70B. Structured prompt: passage + rating rubric ($-2$ strongly liberal to $+2$ strongly conservative). Deterministic decoding (temperature = 0) for reproducibility across five dimensions: word choice, framing, source selection, omission patterns, and tone.

\subsubsection{Bayesian Hierarchical Model}

Three-level hierarchy: dimension $\to$ publisher $\to$ textbook $\to$ passage. Partial pooling provides shrinkage toward grand mean, preventing overfitting to small publishers. MCMC inference: PyMC, 4 chains $\times$ 2,000 draws (1,000 warm-up), NUTS sampler. Convergence: R-hat $< 1.01$ on all parameters; ESS $> 3,000$.

\subsection{Results}

\begin{table}[H]
\centering
\caption{Inter-rater reliability across three LLMs}
\begin{tabular}{lcc}
\toprule
\textbf{Metric} & \textbf{Value} & \textbf{Interpretation} \\
\midrule
Krippendorff's $\alpha$ (3-LLM) & 0.84 & Excellent ($\geq 0.80$ threshold) \\
GPT-4 $\leftrightarrow$ Claude-3 ($r$) & 0.92 & Near-perfect \\
GPT-4 $\leftrightarrow$ Llama-3 ($r$) & 0.89 & Excellent \\
Claude-3 $\leftrightarrow$ Llama-3 ($r$) & 0.87 & Excellent \\
Fleiss' $\kappa$ (3-way) & 0.81 & Substantial to excellent \\
\bottomrule
\end{tabular}
\label{tab:llm_reliability}
\end{table}

\begin{table}[H]
\centering
\caption{Publisher bias estimates from Bayesian hierarchical model}
\begin{tabular}{llccc}
\toprule
\textbf{Rank} & \textbf{Publisher} & \textbf{Posterior Mean} & \textbf{95\% HDI} & \textbf{Classification} \\
\midrule
1 & Publisher C & $-0.48$ & $[-0.62, -0.34]$ & \textbf{Credibly Liberal} \\
2 & Publisher A & $-0.29$ & $[-0.41, -0.17]$ & \textbf{Credibly Liberal} \\
3 & Publisher E & $+0.02$ & $[-0.10, +0.14]$ & Neutral \\
4 & Publisher B & $+0.08$ & $[-0.04, +0.20]$ & Neutral \\
5 & Publisher D & $+0.38$ & $[+0.26, +0.50]$ & \textbf{Credibly Conservative} \\
\bottomrule
\end{tabular}
\label{tab:publisher_bias}
\end{table}

\noindent\textit{Classification rule: credible effect when 95\% HDI excludes zero. Scale: $-2$ (strongly liberal) to $+2$ (strongly conservative).}

\begin{table}[H]
\centering
\caption{MCMC convergence diagnostics}
\begin{tabular}{lccc}
\toprule
\textbf{Parameter Group} & \textbf{R-hat} & \textbf{ESS (bulk)} & \textbf{ESS (tail)} \\
\midrule
Publisher effects & 1.001 & 4,218 & 3,891 \\
Textbook effects & 1.002 & 3,744 & 3,512 \\
Dimension effects & 1.001 & 4,106 & 3,988 \\
Grand mean & 1.000 & 5,201 & 4,873 \\
\bottomrule
\end{tabular}
\label{tab:mcmc}
\end{table}

\subsection{Threats to Validity}

\begin{itemize}[leftmargin=*]
    \item LLMs may share training data with evaluated textbooks, introducing systematic rating bias
    \item Five-publisher corpus may not generalize to independent, state-specific, or digital-native publishers
    \item Political bias scale is ordinal; interval assumptions in Bayesian model are approximations
    \item LLM political calibration is time-sensitive; model updates may alter future ratings without replication
\end{itemize}

%=============================================================================
\section{Cross-Project Technical Strengths}
%=============================================================================

\subsection{Statistical Rigor}

\begin{table}[H]
\centering
\caption{Statistical capabilities by project}
\begin{tabular}{lccc}
\toprule
\textbf{Capability} & \textbf{Project 1} & \textbf{Project 2} & \textbf{Project 3} \\
\midrule
Confidence/HDI intervals & \checkmark (95\% HDI) & \checkmark (95\% CI) & \checkmark (95\% HDI) \\
Cross-validation & \checkmark (10-fold) & \checkmark (10-fold) & N/A \\
Inter-rater reliability & \checkmark ($\alpha = 0.81$) & N/A & \checkmark ($\alpha = 0.84$) \\
Calibration analysis & \checkmark & \checkmark (ECE, Brier) & N/A \\
Bayesian inference & \checkmark (hierarchical) & N/A & \checkmark (hierarchical) \\
Multiple testing adjustment & \checkmark & N/A & \checkmark (Dunn/BH) \\
\bottomrule
\end{tabular}
\label{tab:stats}
\end{table}

\subsection{Responsible AI Controls}

All three projects include:

\begin{itemize}[leftmargin=*]
    \item \textbf{Limitations disclosure}: explicit threats to validity and generalizability caveats
    \item \textbf{Human-in-the-loop}: configurable escalation thresholds and audit logging
    \item \textbf{Governance documentation}: model cards, intended-use statements, misuse risk assessment
    \item \textbf{Uncertainty communication}: all estimates carry quantified uncertainty intervals
\end{itemize}

%=============================================================================
\section{Role Alignment and Competency Map}
%=============================================================================

\begin{table}[H]
\centering
\caption{Competency evidence map}
\begin{tabular}{lp{9cm}}
\toprule
\textbf{Competency} & \textbf{Evidence} \\
\midrule
Supervised ML (classification) & Projects 1 \& 2: 8+ ensemble algorithms, full benchmark \\
Bayesian inference & Projects 1 \& 3: PyMC hierarchical models with MCMC \\
LLM evaluation \& prompting & Projects 1 \& 3: multi-LLM annotation, structured prompts \\
Calibration \& uncertainty & Project 2: ECE, Brier, Platt scaling \\
Statistical hypothesis testing & Projects 2 \& 3: binomial, Kruskal--Wallis, Dunn, Friedman \\
Explainability (SHAP) & Projects 1 \& 2: global + local attribution \\
MLOps \& deployment & Projects 1 \& 2: REST API, latency profiling, drift monitoring \\
Responsible AI & All three: governance docs, limitations, escalation policies \\
\bottomrule
\end{tabular}
\label{tab:competency}
\end{table}

\begin{table}[H]
\centering
\caption{Applicable role alignment}
\begin{tabular}{ll}
\toprule
\textbf{Role} & \textbf{Alignment} \\
\midrule
Data Scientist (Applied ML) & Primary (Projects 1 \& 2) \\
Applied Statistician & Primary (Projects 2 \& 3) \\
GenAI Evaluation Scientist & Primary (Projects 1 \& 3) \\
ML Research Engineer & All three projects \\
AI Safety / Red Team Researcher & Project 1 \\
\bottomrule
\end{tabular}
\label{tab:roles}
\end{table}

%=============================================================================
\section{Reproducibility and Deployment Notes}
%=============================================================================

\subsection{Code and Data Availability}

\begin{table}[H]
\centering
\caption{Repository resources}
\begin{tabular}{ll}
\toprule
\textbf{Resource} & \textbf{Location} \\
\midrule
Repository & \url{github.com/dl1413/Machine-Learning-Research-Engineering-Project-Profile} \\
AI Safety Notebook & \texttt{AI\_Safety\_RedTeam\_Evaluation.ipynb} \\
Breast Cancer Notebook & \texttt{Breast\_Cancer\_Classification\_PUBLICATION.ipynb} \\
Bias Detection Notebook & \texttt{LLM\_Ensemble\_Textbook\_Bias\_Detection.ipynb} \\
PDF Build Pipeline & \texttt{scripts/build\_reports\_pdf.sh} \\
\bottomrule
\end{tabular}
\label{tab:resources}
\end{table}

\subsection{Reproducibility Controls}

All projects share: fixed random seeds (seed = 42 throughout), pinned dependency versions (requirements files), deterministic preprocessing pipelines, and environment specification (Python 3.11, conda/pip lockfiles).

\subsection{PDF Generation}

Publication-ready PDFs are generated via:

\begin{verbatim}
chmod +x scripts/build_reports_pdf.sh
./scripts/build_reports_pdf.sh
\end{verbatim}

Output PDFs are in \texttt{pdf/out/}. See \texttt{PDF\_EXPORT.md} for full build instructions.

%=============================================================================
\section{Conclusions}
%=============================================================================

This portfolio demonstrates end-to-end machine learning research engineering capability across three distinct applied domains:

\begin{enumerate}[leftmargin=*]
    \item \textbf{AI Safety Red-Team Evaluation} establishes a cost-effective, audit-grade framework for scalable LLM harm detection, reducing annotation cost by 340$\times$ while preserving expert-level reliability ($\alpha = 0.81$, ROC-AUC 0.9923).

    \item \textbf{Breast Cancer Classification} delivers clinically viable diagnostic support with exceptional discriminative performance (ROC-AUC 0.9987) and improved probability calibration (ECE reduced from 0.0312 to 0.0089), enabling threshold-optimized operating points for screening versus diagnosis contexts.

    \item \textbf{LLM Ensemble Bias Detection} provides a reproducible, statistically rigorous audit methodology for educational content, identifying credible political framing effects in 3 of 5 publishers with an uncertainty-aware triage protocol for expert review.
\end{enumerate}

Together, these projects reflect competencies in ML systems design, statistical inference, responsible AI governance, and engineering rigor required for senior data science and applied ML research roles in 2026.

%=============================================================================
\begin{thebibliography}{14}

\bibitem{krippendorff2011}
Krippendorff, K. (2011).
\newblock Computing Krippendorff's alpha-reliability.
\newblock \textit{Annals of Applied Statistics}, 5(1), 103--117.

\bibitem{chen2016xgboost}
Chen, T., \& Guestrin, C. (2016).
\newblock XGBoost: A scalable tree boosting system.
\newblock \textit{Proceedings of KDD '16}, 785--794.

\bibitem{ke2017lightgbm}
Ke, G., et al. (2017).
\newblock LightGBM: A highly efficient gradient boosting decision tree.
\newblock \textit{NeurIPS 30}.

\bibitem{gelman2013}
Gelman, A., et al. (2013).
\newblock \textit{Bayesian Data Analysis} (3rd ed.).
\newblock Chapman \& Hall/CRC.

\bibitem{chawla2002smote}
Chawla, N.V., et al. (2002).
\newblock SMOTE: Synthetic minority over-sampling technique.
\newblock \textit{Journal of Machine Learning Research}, 3, 321--357.

\bibitem{lundberg2017shap}
Lundberg, S., \& Lee, S. (2017).
\newblock A unified approach to interpreting model predictions.
\newblock \textit{NeurIPS 30}.

\bibitem{mangasarian1995}
Mangasarian, O.L., et al. (1995).
\newblock Breast cancer diagnosis and prognosis via linear programming.
\newblock \textit{Operations Research}, 43(4), 570--577.

\bibitem{brier1950}
Brier, G.W. (1950).
\newblock Verification of forecasts expressed in terms of probability.
\newblock \textit{Monthly Weather Review}, 78(1), 1--3.

\bibitem{openai2024gpt4}
OpenAI. (2024).
\newblock GPT-4 technical report.
\newblock arXiv:2303.08774.

\bibitem{anthropic2024claude}
Anthropic. (2024).
\newblock Claude 3 model card.
\newblock \textit{Anthropic Technical Report}.

\bibitem{ieee2025}
IEEE. (2025).
\newblock IEEE 2830-2025: Standard for transparent ML documentation.

\bibitem{iso2025}
ISO/IEC. (2025).
\newblock ISO/IEC 23894:2025: Artificial intelligence --- Risk management.

\bibitem{euai2025}
European Commission. (2025).
\newblock EU Artificial Intelligence Act.

\bibitem{niculescu2005}
Niculescu-Mizil, A., \& Caruana, R. (2005).
\newblock Predicting good probabilities with supervised learning.
\newblock \textit{ICML '05}, 625--632.

\end{thebibliography}

\end{document}
